\section{Results}

\subsection{Training Convergence}
\begin{figure}[]
    \centering

    \begin{tikzpicture}
        \begin{axis}[
                grid=major,
                legend entries = {
                    CNN,
                    BiLSTM,
                    MLP,
                    FusionNet,
                    SelfAttention
                },
            ]
            \addplot table {data/model_loss/CNN.dat};
            \addplot table {data/model_loss/BiLSTM.dat};
            \addplot table {data/model_loss/MLP.dat};
            \addplot table {data/model_loss/FusionNet.dat};
            \addplot table {data/model_loss/SelfAttention.dat};
        \end{axis}
    \end{tikzpicture}
    \caption{Training Loss of various models across epochs.}\label{fig:train-loss}
\end{figure}


Figure~\ref{fig:train-loss} presents the training loss curves for all five architectures across five epochs. The MLP exhibits the lowest absolute loss throughout training, declining from 0.0035 to 0.0031, a trajectory that reflects the relative simplicity of its optimization landscape rather than superior generalization capacity. The CNN and BiLSTM follow closely comparable convergence paths, both reaching a terminal loss of approximately 0.0033 after initial values of 0.0042 and 0.0057, respectively. FusionNet and SelfAttention begin training with considerably higher initial losses of 0.0134 and 0.0100, consistent with their greater architectural complexity and the correspondingly more challenging optimization surface. Both descend steeply over the first two epochs; FusionNet converges to 0.0035 by epoch four, while SelfAttention stabilizes at 0.0048, the highest final loss among all evaluated models. The accelerated loss reduction observed in FusionNet following the initial epochs suggests that its dedicated structured feature branch supports convergence once the text encoder begins to stabilize.

\subsection{Multi-Label Classification Performance}
\begin{table}[]
    \caption{Results of the evaluated models across various metrics.}\label{tab:result-metrics}
    \centerline{
    \begin{tabular}{@{}llllllll@{}}
        \toprule
        \textbf{Model} & \textbf{Jaccard}       & \textbf{Precision\_micro} & \textbf{Recall\_micro} & \textbf{F1\_micro} & \textbf{Precision\_macro} & \textbf{Recall\_macro} & \textbf{F1\_macro} \\ \midrule
        MLP            & 0.6887                 & 0.8853                    & 0.7561                 & 0.8156             & 0.0781                    & 0.0675                 & 0.0721             \\
        CNN            & 0.7657                 & 0.9031                    & 0.8342                 & 0.8673             & 0.0858                    & 0.0762                 & 0.0798             \\
        BiLSTM         & 0.7796                 & 0.9218                    & 0.8348                 & 0.8762             & 0.0899                    & 0.0777                 & 0.0821             \\
        SelfAttention  & 0.7684                 & 0.9022                    & 0.8382                 & 0.8691             & 0.0864                    & 0.0753                 & 0.0789             \\
        FusionNet      & 0.7869                 & 0.9252                    & 0.8404                 & 0.8808             & 0.0898                    & 0.0786                 & 0.0828             \\ \midrule
        \textbf{Model} & \textbf{Hamming\_Loss} & \textbf{Subset\_Accuracy} & \textbf{ROC\_AUC}      & \textbf{PRAUC}     & \textbf{Precision@K}      & \textbf{Recall@K}      & \textbf{F1@K}      \\ \midrule
        MLP            & 0.0015                 & 0.4210                    & 0.9937                 & 0.8372             & 0.6855                    & 0.6329                 & 0.5825             \\
        CNN            & 0.0011                 & 0.5603                    & 0.9968                 & 0.8743             & 0.7101                    & 0.6563                 & 0.6042             \\
        BiLSTM         & 0.0010                 & 0.6184                    & 0.9957                 & 0.8743             & 0.7076                    & 0.6555                 & 0.6028             \\
        SelfAttention  & 0.0011                 & 0.5675                    & 0.9970                 & 0.8764             & 0.7114                    & 0.6573                 & 0.6052             \\
        FusionNet      & 0.0010                 & 0.6455                    & 0.9964                 & 0.8753             & 0.7079                    & 0.6561                 & 0.6034             \\ \bottomrule
    \end{tabular}
        }

\end{table}


Table~\ref{tab:result-metrics} reports the full suite of evaluation metrics across all five architectures. FusionNet achieves the strongest overall performance, attaining the highest Jaccard index (0.7869), micro-precision (0.9252), micro-F1 (0.8808), and subset accuracy (0.6455), with a Hamming loss of 0.0010. BiLSTM ranks second across several key metrics, including Jaccard index (0.7796), micro-F1 (0.8762), and subset accuracy (0.6184). SelfAttention and CNN occupy an intermediate performance tier; notably, SelfAttention achieves marginally higher PR-AUC (0.8764) and Precision@$K$ (0.7114) than both CNN and BiLSTM, indicating that attention-based scoring is particularly effective for top-$k$ retrieval tasks. The MLP baseline, while competitive on micro-precision (0.8853), lags substantially on subset accuracy (0.4210) and Jaccard index (0.6887), suggesting that flat feature representations are insufficient for accurate multi-label drug set prediction.

Macro-averaged metrics are uniformly low across all models (macro-F1 $\approx$ 0.07--0.08), an expected outcome given the large and heavily imbalanced drug vocabulary; infrequent drug classes appear rarely during training, yielding poor per-class recall for low-frequency medications.

\subsection{Ranking and Top-$K$ Performance}
All models attain high ROC-AUC scores (${>}0.993$), confirming strong discriminative capacity at the individual label level. PR-AUC scores range from 0.8372 for the MLP to 0.8764 for SelfAttention, with CNN, BiLSTM, and FusionNet clustered closely in the range 0.874--0.875. At $k = 5$, Precision@$K$ spans 0.6855 (MLP) to 0.7114 (SelfAttention), while F1@$K$ ranges from 0.5825 to 0.6052. The narrow spread in top-$k$ metrics among the four stronger architectures suggests that beyond a certain representational threshold, marginal gains in ranking quality diminish, and considerations such as safety filtering and interpretability become the primary differentiating factors.

\subsection{Summary}
FusionNet emerges as the best-performing architecture overall, combining strong multi-label accuracy with the highest subset accuracy and thus representing the most suitable backbone for the proposed medication recommendation framework. Its explicit separation of text and structured feature encoding is well-aligned with the multi-modal nature of outpatient clinical data. SelfAttention offers competitive ranking performance and may be preferred in deployment scenarios where top-$k$ precision constitutes the primary evaluation criterion. The MLP baseline, despite its architectural simplicity, provides a meaningful lower bound and confirms that even shallow models derive substantial benefit from structured feature integration.
